% 한국교원대학교 Beamer 발표 템플릿의 공통 설정 파일입니다.
% Shared settings for the KNUE Beamer presentation template.
%
% 제목, 저자, 참고문헌, 발표자 노트, 자주 쓰는 명령을 이 파일에서 관리합니다.
% This file manages title metadata, bibliography, speaker notes, and reusable helper commands.

% ========== Presentation Information / 발표 정보 ==========

% 아래 명령의 중괄호 안만 자신의 발표 정보로 바꿉니다.
% Replace only the text inside braces with your own presentation information.

% 학위과정 구분입니다. 논문 템플릿처럼 석사 또는 박사 중 하나로 바꿉니다.
% Degree program selector. As in the thesis template, choose either 석사 or 박사.
\providecommand{\degreeProgramKor}{박사} % 석사 또는 박사

% 발표 종류 구분입니다. proposal은 논문 프로포절, result는 논문 결과발표입니다.
% Presentation-stage selector. Use proposal for proposal talks and result for result/defense talks.
\providecommand{\presentationStage}{proposal} % proposal 또는 result

% 선택값을 preamble에서 한 번만 판정합니다. 한글 선택값은 토큰을 문자열로 바꿔 비교해야 안정적입니다.
% Evaluate selector values once in the preamble. Korean selector values are compared as strings for stability.
\usepackage{pdftexcmds}
\makeatletter
\newif\ifKNUEMaster
\newif\ifKNUEResult
\edef\KNUE@mastervalue{\detokenize{석사}}
\edef\KNUE@resultvalue{\detokenize{result}}
\edef\KNUE@degreevalue{\detokenize\expandafter{\degreeProgramKor}}
\edef\KNUE@stagevalue{\detokenize\expandafter{\presentationStage}}
\ifnum\pdf@strcmp{\KNUE@degreevalue}{\KNUE@mastervalue}=0
  \KNUEMastertrue
\else
  \KNUEMasterfalse
\fi
\ifnum\pdf@strcmp{\KNUE@stagevalue}{\KNUE@resultvalue}=0
  \KNUEResulttrue
\else
  \KNUEResultfalse
\fi
\makeatother

% 학위과정에 따라 국문 학위논문 문구를 만듭니다.
% Build the Korean thesis/dissertation label from the degree program.
\newcommand{\BeamerDegreeWorkKor}{\degreeProgramKor{}학위논문}

% 학위과정에 따라 영문 학위논문 문구를 미리 정합니다.
% Predefine the English thesis/dissertation label from the degree program.
\ifKNUEMaster
  \newcommand{\BeamerDegreeWorkEng}{Master's Thesis}%
\else
  \newcommand{\BeamerDegreeWorkEng}{Doctoral Dissertation}%
\fi

% 발표 종류에 따라 국문/영문 발표 문구를 미리 정합니다.
% Predefine Korean/English presentation-stage labels from the stage selector.
\ifKNUEResult
  \newcommand{\BeamerStageKor}{결과발표}%
  \newcommand{\BeamerStageEng}{Results Presentation}%
\else
  \newcommand{\BeamerStageKor}{연구계획 발표}%
  \newcommand{\BeamerStageEng}{Proposal Presentation}%
\fi

% 발표 종류에 따라 두 가지 내용 중 하나를 골라 출력합니다.
% Choose between two texts according to the selected presentation stage.
\newcommand{\KNUEStageText}[2]{%
  \ifKNUEResult #2\else #1\fi%
}

% 발표 종류에 따라 결과 섹션의 제목 문구를 미리 정합니다.
% Predefine result-section labels according to the selected presentation stage.
\ifKNUEResult
  \newcommand{\BeamerResultsSectionTitle}{연구 결과}%
  \newcommand{\BeamerResultsSectionTitleEng}{Research Results}%
  \newcommand{\BeamerResultsFirstSubsection}{주요 결과}%
  \newcommand{\BeamerResultsFirstFrame}{주요 연구 결과}%
  \newcommand{\BeamerResultsFirstBlock}{연구를 통해 확인한 결과}%
  \newcommand{\BeamerResultsSecondSubsection}{결과 해석}%
  \newcommand{\BeamerResultsSecondFrame}{결과 해석}%
  \newcommand{\BeamerResultsThirdSubsection}{논의와 시사점}%
  \newcommand{\BeamerResultsThirdFrame}{논의와 시사점}%
  \newcommand{\BeamerSummaryFrame}{연구 결과 요약}%
  \newcommand{\BeamerQASubsection}{질의 및 논의}%
  \newcommand{\BeamerQAFrame}{질의 및 논의 사항}%
  \newcommand{\BeamerQABlock}{심사위원과 논의할 부분}%
  \newcommand{\BeamerDefaultShortTitle}{Results}%
\else
  \newcommand{\BeamerResultsSectionTitle}{예상 결과}%
  \newcommand{\BeamerResultsSectionTitleEng}{Expected Results}%
  \newcommand{\BeamerResultsFirstSubsection}{예상 산출물}%
  \newcommand{\BeamerResultsFirstFrame}{예상 산출물}%
  \newcommand{\BeamerResultsFirstBlock}{연구를 통해 산출할 결과}%
  \newcommand{\BeamerResultsSecondSubsection}{연구 일정}%
  \newcommand{\BeamerResultsSecondFrame}{연구 일정}%
  \newcommand{\BeamerResultsThirdSubsection}{기대 효과}%
  \newcommand{\BeamerResultsThirdFrame}{기대 효과}%
  \newcommand{\BeamerSummaryFrame}{연구계획 요약}%
  \newcommand{\BeamerQASubsection}{검토 요청}%
  \newcommand{\BeamerQAFrame}{검토 요청 사항}%
  \newcommand{\BeamerQABlock}{심사위원께 검토를 요청할 부분}%
  \newcommand{\BeamerDefaultShortTitle}{Proposal}%
\fi

% 국문 발표 제목입니다. 표지에서 가장 크게 출력됩니다.
% Korean presentation title. This is shown most prominently on the title slide.
\providecommand{\BeamerTitleKor}{학위논문 발표 제목}

% 영문 발표 제목입니다. 국문 제목 아래에 작은 글씨로 출력됩니다.
% English presentation title. This appears below the Korean title in smaller type.
\providecommand{\BeamerTitleEng}{Thesis or Dissertation Presentation Title}

% 발표의 성격을 짧게 설명하는 국문 부제입니다.
% Korean subtitle that describes the presentation type.
\providecommand{\BeamerSubtitleKor}{\BeamerDegreeWorkKor{} \BeamerStageKor}

% 발표의 성격을 짧게 설명하는 영문 부제입니다.
% English subtitle that describes the presentation type.
\providecommand{\BeamerSubtitleEng}{\BeamerDegreeWorkEng{} \BeamerStageEng}

% PDF 북마크와 사이드바에서 사용할 짧은 제목입니다.
% Short title used in PDF bookmarks and the sidebar.
\providecommand{\BeamerShortTitle}{\BeamerDefaultShortTitle}

% 발표자 이름입니다.
% Presenter name.
\providecommand{\BeamerAuthor}{홍 길 동}

% 지도교수 이름입니다. 필요하면 "지도교수:" 문구까지 함께 수정합니다.
% Advisor name. Edit the "지도교수:" label too if needed.
\providecommand{\BeamerAdvisor}{지도교수: 성 춘 향}

% 소속 학과, 전공, 연구실 정보를 적습니다.
% Department, major, lab, or institutional affiliation.
\providecommand{\BeamerInstitute}{한국교원대학교 지구과학교육전공}

% 발표일입니다. 기본값 \today는 빌드한 날짜를 자동으로 넣습니다.
% Presentation date. The default \today inserts the build date automatically.
\providecommand{\BeamerDate}{\today}

% 이전 연구계획서 템플릿 명령을 계속 사용할 수 있게 남겨 둔 호환 별칭입니다.
% Compatibility aliases for documents created from the earlier proposal template.
\providecommand{\ProposalTitleKor}{\BeamerTitleKor}
\providecommand{\ProposalTitleEng}{\BeamerTitleEng}
\providecommand{\ProposalSubtitle}{\BeamerSubtitleKor}
\providecommand{\ProposalAuthor}{\BeamerAuthor}
\providecommand{\ProposalAdvisor}{\BeamerAdvisor}
\providecommand{\ProposalInstitute}{\BeamerInstitute}
\providecommand{\ProposalDate}{\BeamerDate}

% ========== Packages / 패키지 ==========

% multicol은 긴 목차를 두 단으로 나눌 때 사용합니다.
% multicol is used to split long table-of-contents slides into two columns.
\usepackage{multicol}

% csquotes는 biblatex-apa와 함께 인용부호 처리를 안정적으로 해 줍니다.
% csquotes works with biblatex-apa to handle quotation marks consistently.
\usepackage{csquotes}

% APA 7판 형식 참고문헌을 biber로 처리합니다.
% Use biber to process APA 7th edition references.
\usepackage[
  backend=biber,
  style=apa,
  natbib=true,
  maxbibnames=20,
  minbibnames=20,
  maxcitenames=2,
  mincitenames=1,
  uniquelist=false,
  dashed=false,
  date=year,
  urldate=iso,
  dateera=astronomical,
  datezeros=true,
  seconds=true
]{biblatex}

% 논문 템플릿과 동일하게 biblatex-apa의 APA 7판 지역화 규칙을 사용합니다.
% Use the same APA 7 localization rules as the thesis template.
\DeclareLanguageMapping{english}{english-apa}
\DeclareLanguageMapping{korean}{english-apa}

% 참고문헌 파일 경로입니다. 메인 파일에서 다른 경로를 지정하면 그 값을 우선 사용합니다.
% Bibliography path. If the main file defines another path, that value takes precedence.
\providecommand{\beamerbib}{sub/references.bib}
\providecommand{\proposalbib}{\beamerbib}

% 실제 참고문헌 데이터베이스를 biblatex에 등록합니다.
% Register the bibliography databases with biblatex.
\addbibresource{\beamerbib}

% pgfpages는 발표자 노트를 두 번째 화면에 배치할 때 필요합니다.
% pgfpages is needed when speaker notes are placed on a second screen.
\usepackage{pgfpages}

% ========== Bibliography and Citation Settings / 참고문헌과 인용 설정 ==========

% references.bib의 langid=korean 문헌을 먼저, langid=english 문헌을 나중에 정렬합니다.
% Sort langid=korean entries first and langid=english entries later.
\DeclareSourcemap{
  \maps[datatype=bibtex]{
    \map{
      \step[fieldsource=langid, match=\regexp{korean}, final]
      \step[fieldset=presort, fieldvalue={a}]
    }
    \map{
      \step[fieldsource=langid, match=\regexp{english}, final]
      \step[fieldset=presort, fieldvalue={z}]
    }
  }
}

% 참고문헌 목록에서 한국어 문헌은 저자명을 가운데점으로 연결하고, 영어 문헌은 APA 7판 기본 연결 방식을 사용합니다.
% Korean bibliography entries join author names with centered dots; English entries keep APA 7 delimiters.
\DeclareDelimFormat[bib,biblist]{multinamedelim}{%
  \iffieldequalstr{langid}{korean}
  {\textperiodcentered}
  {\addcomma\space}%
}
\DeclareDelimFormat[bib,biblist]{finalnamedelim}{%
  \iffieldequalstr{langid}{korean}
  {\textperiodcentered}
  {\ifnumgreater{\value{liststop}}{2}{\finalandcomma}{}\addspace\&\space}%
}
\DeclareDelimFormat[bib,biblist]{nameyeardelim}{%
  \iffieldequalstr{langid}{korean}
  {}
  {\newunitpunct}%
}
\DeclareDelimAlias[bib,biblist]{nonameyeardelim}{nameyeardelim}

% 한국어 학술지명과 권은 굵게, 영어 학술지명과 권은 APA 기본 기울임꼴로 표시합니다.
% Korean journal titles and volumes are bold; English ones remain APA-style italics.
\DeclareFieldFormat[article]{journaltitle}{%
  \iffieldequalstr{langid}{korean}
  {\mkbibbold{#1}}
  {\mkbibemph{#1}}%
}
\DeclareFieldFormat[article]{volume}{%
  \iffieldequalstr{langid}{korean}
  {\mkbibbold{#1}}
  {\mkbibemph{#1}}%
}

% DOI와 URL은 클릭 가능한 주소 형태로 표시합니다.
% Display DOI and URL fields as clickable address-style links.
\DeclareFieldFormat{doi}{%
  \ifhyperref{\href{https://doi.org/#1}{\nolinkurl{https://doi.org/#1}}}{\nolinkurl{https://doi.org/#1}}%
  \nopunct%
}
\DeclareFieldFormat{url}{\url{#1}}

% 서술형 인용에서 저자명과 연도 사이 공백을 줄입니다.
% Remove the space between author names and year in textual citations.
\DeclareDelimFormat[textcite]{nameyeardelim}{}

% Beamer 참고문헌은 본문보다 작지만 읽을 수 있도록 scriptsize로 설정합니다.
% Use scriptsize for Beamer bibliography entries so they remain readable.
\renewcommand*{\bibfont}{\scriptsize}

% 참고문헌의 내어쓰기 폭을 줄여 좁은 슬라이드에서도 줄바꿈이 덜 부담스럽게 보이게 합니다.
% Use a shorter hanging indent so wrapped bibliography lines fit Beamer slides better.
\setlength{\bibhang}{0.8em}

% 참고문헌 항목 사이의 간격을 줄여 한 슬라이드에 더 많이 들어가게 합니다.
% Reduce spacing between bibliography entries so more items fit on slides.
\setlength\bibitemsep{0pt}
\setlength\bibnamesep{0pt}

% 참고문헌 아이콘 대신 텍스트형 항목 표시를 사용합니다.
% Use text-style bibliography items instead of icon bullets.
\setbeamertemplate{bibliography item}[text]

% 참고문헌 제목을 한국어로 보이게 합니다.
% Display bibliography headings in Korean.
\DefineBibliographyStrings{english}{
  references = {참고문헌},
  bibliography = {참고문헌},
}

% ---------- 한국어 문헌 본문 인용 명령 / Korean citation commands ----------
% 한국어 문헌은 \ktextcite, \kparencite를 사용하고 영어 문헌은 \textcite, \parencite를 사용합니다.
% Use \ktextcite and \kparencite for Korean sources; use \textcite and \parencite for English sources.
% \kparencite와 \parencite는 괄호를 직접 출력하지 않습니다.
% \kparencite and \parencite do not print outer parentheses; write them in the slide text when needed.

\newcommand{\KoreanCiteAuthorDelim}{\textperiodcentered}
\newcommand{\UseKoreanCiteDelims}{%
  \DeclareDelimFormat{multinamedelim}{\KoreanCiteAuthorDelim}%
  \DeclareDelimFormat[textcite]{multinamedelim}{\KoreanCiteAuthorDelim}%
  \DeclareDelimFormat[parencite]{multinamedelim}{\KoreanCiteAuthorDelim}%
  \DeclareDelimFormat{finalnamedelim}{\KoreanCiteAuthorDelim}%
  \DeclareDelimFormat[textcite]{finalnamedelim}{\KoreanCiteAuthorDelim}%
  \DeclareDelimFormat[parencite]{finalnamedelim}{\KoreanCiteAuthorDelim}%
  \DeclareDelimFormat[textcite]{nameyeardelim}{}%
  \DeclareDelimFormat{andothersdelim}{\addspace}%
  \DeclareDelimFormat[textcite]{andothersdelim}{\addspace}%
  \DeclareDelimFormat[parencite]{andothersdelim}{\addspace}%
  \csdef{abx@lstr@andothers}{외}%
  \csdef{abx@sstr@andothers}{외}%
}
\let\KNUEorigparencite\parencite
\let\parencite\cite
\NewDocumentCommand{\ktextcite}{o o m}{%
  \begingroup
  \UseKoreanCiteDelims
  \IfNoValueTF{#1}{%
    \textcite{#3}%
  }{%
    \IfNoValueTF{#2}{\textcite[#1]{#3}}{\textcite[#1][#2]{#3}}%
  }%
  \endgroup
}
\NewDocumentCommand{\kparencite}{o o m}{%
  \begingroup
  \UseKoreanCiteDelims
  \IfNoValueTF{#1}{%
    \cite{#3}%
  }{%
    \IfNoValueTF{#2}{\cite[#1]{#3}}{\cite[#1][#2]{#3}}%
  }%
  \endgroup
}

% ========== Table of Contents / 목차 설정 ==========

% 목차에서 section과 subsection 앞에 붙일 점 표식입니다.
% Bullet marker used before section and subsection entries in the table of contents.
\newcommand{\knueTocBullet}{\raisebox{0.15ex}{\large\textbullet}}

% 목차에서 section은 subsection보다 한 단계 크게 표시합니다.
% Sections are displayed one size larger than subsections in ToC slides.
\setbeamerfont{section in toc}{size=\small}
\setbeamerfont{subsection in toc}{size=\scriptsize}

% section 목차 항목의 줄바꿈과 들여쓰기를 안정화합니다.
% Stabilize wrapping and indentation for section entries in the ToC.
\setbeamertemplate{section in toc}{%
  \vspace{0.35em}%
  \usebeamerfont{section in toc}%
  \usebeamercolor[fg]{section in toc}%
  \leavevmode\parindent0pt%
  \hangindent=1.4em\hangafter=1%
  \makebox[1.2em][l]{\knueTocBullet}\inserttocsection\par%
}

% subsection 목차 항목은 section보다 안쪽에 들여쓰기합니다.
% Subsection entries are indented relative to section entries.
\setbeamertemplate{subsection in toc}{%
  \vspace{0.15em}%
  \usebeamerfont{subsection in toc}%
  \usebeamercolor[fg]{subsection in toc}%
  \leavevmode\parindent0pt%
  \leftskip=1.4em%
  \hangindent=1.4em\hangafter=1%
  \makebox[1.2em][l]{\knueTocBullet}\inserttocsubsection\par%
}

% ========== Presenter Notes / 발표자 노트 ==========

% overlay에서 아직 드러나지 않은 항목을 희미하게 보이도록 설정합니다.
% Show not-yet-uncovered overlay items transparently.
\setbeamercovered{transparent}

% 노트 페이지의 기본 글자 크기입니다.
% Default font size for note pages.
\setbeamerfont{note page}{size=\footnotesize}

% 발표자 노트를 두 번째 화면 오른쪽에 출력하는 명령입니다.
% Command that shows speaker notes on the right side of a second screen.
\newcommand{\enablenotes}{\setbeameroption{show notes on second screen=right}}

% 발표자 노트를 PDF 출력에서 숨기는 명령입니다.
% Command that hides speaker notes in the PDF output.
\newcommand{\disablenotes}{\setbeameroption{hide notes}}

% 각 슬라이드에 발표자 노트를 쉽게 추가하기 위한 짧은 명령입니다.
% Convenience command for adding speaker notes to each slide.
\newcommand{\slidenote}[1]{\note{\footnotesize #1}}

% 기본값은 청중용 PDF입니다. 노트가 보이지 않습니다.
% Default output is an audience PDF. Notes are hidden.
\disablenotes

% ========== Convenience Macros / 편의 명령 ==========

% 그림, 표, 자료의 출처를 작은 글씨로 표시합니다.
% Display source notes for figures, tables, or materials in tiny text.
\newcommand{\source}[1]{%
  \vspace{0.15em}%
  {\tiny 출처: #1}%
}

% 짧은 직접 인용문과 문헌 인용을 함께 넣을 때 사용합니다.
% Use this for a short quoted phrase followed by a citation.
\newcommand{\qcite}[2]{\enquote{#1} (\parencite{#2})}

% itemize/enumerate의 간격을 줄여 발표 슬라이드에 맞게 만듭니다.
% Tighten itemize/enumerate spacing for presentation slides.
\newcommand{\tightitems}{%
  \setlength{\itemsep}{2pt}%
  \setlength{\parskip}{0pt}%
}
